\documentclass[10pt,a4paper]{amsart}
\usepackage[margin=19mm]{geometry}
\usepackage{amsmath,amssymb,mathtools}
\usepackage[scr=rsfs,cal=euler]{mathalfa}
\usepackage{xfrac}
\usepackage[unicode,hidelinks,bookmarks=false]{hyperref}
\newcommand{\ie}{i.e.\ }
\newcommand{\cf}{cf.\ }
\newcommand{\resp}{respectively\ }
\newcommand{\Subsection}{\subsection}
\providecommand{\nobreakdash}{\nobreak\hbox{-}}
\setlength{\emergencystretch}{2em}
\multlinegap0pt
\usepackage{bm}
\usepackage{bbm}
\usepackage{dsfont}
\usepackage{xspace}
\usepackage{cancel}
\usepackage{mathtools}
\usepackage{amsthm}
\makeatletter
\@ifundefined{lemma}{\newtheorem{lemma}{Lemma}}{}
\makeatother
\title[Identifiability and Estimation in Continuous Lyapunov Models]{Identifiability and Estimation in Continuous Lyapunov Models}
\author{Cecilie Olesen Recke, Niels Richard Hansen}
\date{Source: arXiv:2603.17142v1 (CC BY 4.0)}
\begin{document}
\maketitle

\noindent\emph{Source and licence.} Identifiability and Estimation in Continuous Lyapunov Models. Version arXiv:2603.17142v1 (CC BY 4.0), distributed under the Creative Commons Attribution 4.0 International licence, as recorded in the arXiv licence metadata for this exact version and shown on the abstract page https://arxiv.org/abs/2603.17142v1. Full rights notice: licenses/arxiv-2603.17142-CC-BY-4.0.txt.\par\smallskip

\noindent\textit{Editorial scope.} Counted unit: the Lemma whose source label is \texttt{lem::forsum} in the article (source0.tex lines 1511--1519, source label \texttt{lem::forsum}), delivered together with the article's own complete original proof of it (source0.tex lines 1520--1569, one \texttt{proof} environment of 50 lines). No mathematics is written, repaired or re-derived: statement and proof are the article's own text line for line. Required context retained uncounted: none. Every definition, display and label of those lines is retained unchanged. Omitted from this package and not redistributed: 1--1510, 1570--2426. Nothing inside a retained excerpt is omitted or altered: every retained body is delivered exactly as the article's lines are written, so no omission receipt of any kind accompanies this package. Citations: the retained excerpts carry no citation command at all, so no bibliography entry of the article is transcribed and no reference is redistributed. Reference adaptation: none was needed and none was made; every reference inside the delivered unit resolves inside it, so no unresolved reference or citation is produced. Printed slips kept literal: no printed word, symbol, hypothesis or number of the counted statement or of its proof was altered. Layout and class: the article's own class is \texttt{biometrika} with packages including amsmath, pgf, tikz, tikz-cd, graphicx, float, hyperref, enumitem, nicematrix, newtxtext, newtxmath, mathtools, algorithm2e; the wrapper is the lane's pinned \texttt{amsart} editorial wrapper at 10pt on A4, which restates the theorem-like environments the retained text needs and adds the article's own macro definitions that the retained text needs (none), with extra wrapper packages (bm, bbm, dsfont, xspace, cancel, mathtools). The delivered numbering is the unit's own. Assets: no retained span contains a figure environment, a graphics-inclusion command, a PDF-page inclusion, a table or a TikZ picture; no image file is copied into this package, the package declares no asset dependency and no image file is redistributed with it. Licence: version arXiv:2603.17142v1 of this article is distributed under the Creative Commons Attribution 4.0 International licence, as recorded in the arXiv licence metadata for this exact version and shown on the abstract page https://arxiv.org/abs/2603.17142v1. A full rights notice is delivered at licenses/arxiv-2603.17142-CC-BY-4.0.txt. Characters: every retained excerpt is pure ASCII, so the delivered engine, XeLaTeX with its default Latin Modern text font, reports no missing character.
\begin{lemma}
\label{lem::forsum}
    Let $d,q$ be non-negative integers such that $q \leq d-1$ and $r \geq 3$ an integer, then 
    \begin{equation}
    \label{eq::idforsumfordeterminant}
            \sum_{i = 0}^{d-1} \left(\frac{r}{2} \right)^{d-1-i} (-1)^{i} c(d,q,i)  = \left(\frac{r}{2} -1 \right)^{d-1} (-1)^{q},
    \end{equation}
where the coefficients are $$c(d,q,i) = \sum_{j = 0}^{i } (r-1)^{j} \binom{q}{j} \binom{d-1-q}{i-j}.$$    
\end{lemma}

\begin{proof}[]
We first prove the following recursive identity for the coefficients
\begin{equation*}
    c(d,q,i) = c(d, q-1, i) + (r-2)c(d-1, q-1, i-1). 
\end{equation*}

To prove this, note that $c(d,q,i)$ is the coefficient in front of $x^{i}$ if you expand the following polynomial in $x$ using the binomial theorem, $(1+(r-1)x)^{q} (1+x)^{d-1-q}$. Thus, $c(d,q-1,i)$ is the coefficient in front of $x^{i}$ in $(1+(r-1)x)^{q-1} (1+x)^{d-1-(q-1)}$ and $c(d-1, q-1, i-1)$ is the coefficient in front of $x^{i-1}$ in $(1+(r-1)x)^{q-1} (1+x)^{d-2-(q-1)}$. Hence, $c(d, q-1, i) + c(d-1, q-1, i-1)$ will be the coefficient in front of $x^{i}$ in 
\begin{align*}
    &(1+(r-1)x)^{q-1} (1+x)^{d-1-(q-1)} + (r-2)x (1+(r-1)x)^{q-1} (1+x)^{d-2-(q-1)}  \\
    &=  (1+(r-1)x)^{q-1}(1+x)^{d-2-(q-1)} (1+x +(r-2)x) \\
    &=  (1+(r-1)x)^{q} (1+x)^{d-1-q}. 
\end{align*}
Thus, the recursion formula is true. 

We now obtain the result about the summation by induction on $(d,q)$, where we induct over this set by what the sum $d + q$ is. 

The base case is $d = 2$ and $q = 0$. The proof is the same for any $d$, so we just prove the formula for any $d$ and $q = 0$. When $q$ is 0 we obtain 
\begin{align*}
    c(d,0,i) = \sum_{j = 0}^{i } (r-1)^{j} \binom{0}{j} \binom{d-1}{i-j} = \binom{d-1}{i},
\end{align*}
since the only non-zero contribution to the sum above is when $j = 0$. Inserting this into the left hand side of equation \eqref{eq::idforsumfordeterminant} we can prove the base case by application of the binomial formula
\begin{align*}
    \sum_{i = 0}^{d-1} \left(\frac{r}{2} \right)^{d-1-i} (-1)^{i} c(d,0,i) = \sum_{i = 0}^{d-1} \left(\frac{r}{2} \right)^{d-1-i} (-1)^{i} \binom{d-1}{i} = \left( \frac{r}{2} -1 \right)^{d-1}.  
\end{align*}

%\textcolor{red}{todo: Can also do p = 2 and q =1 and actually also for any p and q = 0, figure out which is the nicest to put in.}

We now consider the induction step, so consider the lefthand side in equation \eqref{eq::idforsumfordeterminant} and apply the recursion formula to obtain
\begin{align*}
    &\sum_{i = 0}^{d-1} \left(\frac{r}{2} \right)^{d-1-i} (-1)^{i} c(d,q,i) \\ &= 
    \sum_{i = 0}^{d-1} \left(\frac{r}{2} \right)^{d-1-i} (-1)^{i} c(d,q-1,i) +
    (r-2)\sum_{i = 0}^{d-1} \left(\frac{r}{2} \right)^{d-1-i} (-1)^{i} c(d-1,q-1,i-1).  
\end{align*}
By the induction hypothesis the left sum is equal to $\left( r/2 - 1 \right)^{d-1} (-1)^{q-1}$. The rightmost sum can be written as
\begin{align*}
    &\sum_{i = 0}^{d-1} \left(\frac{r}{2} \right)^{d-1-i} (-1)^{i} c(d-1,q-1,i-1) \\ &= - \sum_{i = 0}^{d-1} \left(\frac{r}{2} \right)^{(d-2)-(i-1)} (-1)^{i-1} c(d-1,q-1,i-1) \\
    &= - \sum_{i = -1}^{d-2} \left(\frac{r}{2} \right)^{(d-2)-i} (-1)^{i} c(d-1,q-1,i) \\
    &= - \sum_{i = 0}^{d-2} \left(\frac{r}{2} \right)^{(d-2)-i} (-1)^{i} c(d-1,q-1,i) 
    - \left(\frac{r}{2} \right)^{d-1} c(d-1, q-1, -1) \\
    &= - \left(\frac{r}{2} -1 \right)^{d-2} (-1)^{q-1},
\end{align*}
by using the induction hypothesis and that $c(d-1, q-1, -1) = 0$. Inserting the expressions for the two sums we obtain
\begin{align*}
    \sum_{i = 0}^{d-1} \left(\frac{r}{2} \right)^{d-1-i} (-1)^{i} c(d,q,i) &= \left(\frac{r}{2} -1 \right)^{d-1} (-1)^{q-1} - (r-2)\left(\frac{r}{2} -1 \right)^{d-2} (-1)^{q-1} \\
    &=  \left(\frac{r}{2} -1 \right)^{d-1} (-1)^{q} \left( -1 + \frac{r-2}{r/2-1} \right) \\
    &=  \left(\frac{r}{2} -1 \right)^{d-1} (-1)^{q} \left(\frac{-r/2 + 1 + r-2}{r/2-1} \right) = 
    \left(\frac{r}{2} -1 \right)^{d-1} (-1)^{q},
\end{align*}
as we needed to show. 
\end{proof}

\end{document}
